\section{[INTERNAL] Co-author Comment Tracker}
\label{sec:appendix-internal-tracking}
%% ====================================================================
%% INTERNAL appendix --- visible only when compiling draft.tex.
%% Contents migrated from review/coauthor-comments-pass.md.
%% Do NOT include in submission.
%% ====================================================================

% Pandoc-generated commands: provide fallback definitions.
\providecommand{\tightlist}{\setlength{\itemsep}{0pt}\setlength{\parskip}{0pt}}
\providecommand{\real}[1]{#1}
\providecommand{\pandocbounded}[1]{#1}

\begin{center}
\fbox{\parbox{0.92\textwidth}{\centering\small
\textbf{INTERNAL --- not part of the submission.}\\
This appendix tracks the per-comment response pass on co-author
review notes and is included in the \texttt{draft} build only
(\texttt{draft.tex} sets \texttt{\textbackslash draftmode}).
The \texttt{paper.tex} submission build excludes this section.}}
\end{center}

\bigskip

\textbf{Date:} 2026-04-25 (originally 2026-04-24; refreshed after the Method/Algorithm refactor) \textbf{Purpose:} Systematic review of all coauthor todo-notes and direct edits from the Overleaf (\texttt{master}) branch, with a decision and response for each. This document is the basis for merging the Overleaf edits into the main editorial branch.

\textbf{Legend for the \texttt{Status} column:} - {\color{green!60!black}\faCheckCircle} accepted and applied - {\color{blue!60!black}\faSync} partially applied (needs further discussion) - {\color{orange!80!black}\faPauseCircle} deferred (noted but not yet applied) - {\color{red!70!black}\faTimesCircle} not applied (explanation in response)

\subsection{Status update --- 2026-04-25}\label{status-update-2026-04-25}

Major changes since the original 2026-04-24 pass:

\begin{itemize}
\tightlist
\item
  The Method/Algorithm refactor (commit \texttt{5f67fc8}) renamed \texttt{MLIPAlgorithm} to \texttt{MLIPMethod}, made it \textbf{disjoint} from \texttt{mls:Algorithm} instead of a subclass, and introduced \texttt{hasTrainingAlgorithm} as the single alignment point with ML-Schema. This resolves comment \textbf{4.5} (Yuji: \texttt{mls:Algorithm} alignment mismatch).
\item
  A consolidated Alignment subsection in §4 backed by Table 2 (22 alignment axioms, L1--L22) and Fig. 3 (alignment graph) replaces the three per-module Alignment paragraphs (commit \texttt{2b8e0fc}). The previously unimplemented MDO/CMSO subsumptions are now in the ontology source.
\item
  A \texttt{trainedUsing} shortcut was added on \texttt{TrainedModel}, and \texttt{hasHyperparameterSetting} was relocated from \texttt{BenchmarkResult} to \texttt{MLIPRun} (commit \texttt{2b8e0fc}).
\end{itemize}

The summary table at the end of this document has been updated to reflect both the original per-item resolutions and the new work.

\begin{center}\rule{0.5\linewidth}{0.5pt}\end{center}

\subsection{Section 1 --- Introduction}\label{section-1-introduction}

\subsubsection{1.1 Pranav: citation ordering and missing algorithms}\label{pranav-citation-ordering-and-missing-algorithms}

\begin{quote}
\textbf{Pranav:} HDNNP, GAP might be necessary they lay down earlier foundation. ACE should come first? Why citation are not auto sorted?
\end{quote}

\begin{quote}
\textbf{Jong Hyun (response):} HDNNP and GAP were added, and ACE comes first accordingly. The citations are sorted in a different way for this proceeding. They are sorted alphabetically by the authors' last names.
\end{quote}

\begin{itemize}
\tightlist
\item
  \textbf{Decision:} Accept. HDNNP and GAP are the foundational references and belong in the list.
\item
  \textbf{Daniel's note:} Accepted on 2026-04-24.
\item
  \textbf{Status:} {\color{green!60!black}\faCheckCircle} \textbf{Applied} (commit \texttt{56797e7}). Algorithm list is now: HDNNP, GAP, ACE, MACE, NequIP, M3GNet, MTP.
\end{itemize}

\subsubsection{1.2 Xi: expand ``OWL''}\label{xi-expand-owl}

\begin{quote}
\textbf{Xi:} What is ``OWL''? Should it be written explicitly?
\end{quote}

\begin{quote}
\textbf{Jong Hyun (response):} OWL is expanded here, and the phrase was changed.
\end{quote}

\begin{itemize}
\tightlist
\item
  \textbf{Decision:} Expand OWL only at first occurrence in Section 1 body, not in the abstract. Daniel notes that the abstract is often read independently (indices, conference programs); for the Semantic Web audience, ``OWL'' is well-known; for materials scientists, they will encounter the expansion in the introduction body.
\item
  \textbf{Daniel's note:} Accepted on 2026-04-24 --- body-only expansion.
\item
  \textbf{Status:} {\color{green!60!black}\faCheckCircle} \textbf{Applied} (commit \texttt{a4ba380}). Section 1 first mention now reads ``an ontology expressed in the Web Ontology Language (OWL,2\textasciitilde DL) that fills this gap.''
\end{itemize}

\subsubsection{1.3 Pranav: expand ``RDF''}\label{pranav-expand-rdf}

\begin{quote}
\textbf{Pranav:} may be give full form of RDF as this different than radial distribution function
\end{quote}

\begin{quote}
\textbf{Jong Hyun (response):} RDF was expanded accordingly.
\end{quote}

\begin{itemize}
\tightlist
\item
  \textbf{Decision:} Expand RDF at its first occurrence in the body (Section 2, where RDF is paired with OWL in the discussion of OPTIMADE limitations). Also add W3C spec citations for both RDF and OWL, retroactively adding the OWL citation at its first mention in Section 1.
\item
  \textbf{Daniel's note:} Accepted on 2026-04-24 --- expand and cite.
\item
  \textbf{Status:} {\color{green!60!black}\faCheckCircle} \textbf{Applied} (commit \texttt{3273244}).

  \begin{itemize}
  \tightlist
  \item
    Section 2 now reads: ``federation mechanisms that the Resource Description Framework (RDF)\textasciitilde{}\cite{rdf11concepts} and OWL\textasciitilde{}\cite{owl2overview} provide.''
  \item
    Section 1 first mention of OWL now cites \texttt{owl2overview}.
  \item
    \texttt{references.bib} gains \texttt{rdf11concepts} and \texttt{owl2overview}.
  \end{itemize}
\end{itemize}

\begin{center}\rule{0.5\linewidth}{0.5pt}\end{center}

\subsection{Section 2 --- Related Work}\label{section-2-related-work}

\subsubsection{2.1 Yuji: FAIR citation}\label{yuji-fair-citation}

\begin{quote}
\textbf{Yuji:} Needs citations. https://doi.org/10.1038/sdata.2016.18 https://doi.org/10.1038/s41597-019-0009-6
\end{quote}

\begin{quote}
\textbf{Jong Hyun (response):} The citations were added accordingly.
\end{quote}

\begin{itemize}
\tightlist
\item
  \textbf{Decision:} Accept. The two DOIs are the original FAIR paper and the 2019 addendum.
\item
  \textbf{Daniel's note:} Accepted on 2026-04-24.
\item
  \textbf{Status:} {\color{green!60!black}\faCheckCircle} \textbf{Applied} (commit \texttt{915fb89}). First FAIR mention in Section 2 now cites \texttt{wilkinson2016fair} and \texttt{Wilkinson2019appendum}.
\end{itemize}

\subsubsection{2.2 Pranav: expand ``DFT'' on first use}\label{pranav-expand-dft-on-first-use}

\begin{quote}
\textbf{Pranav:} better say full form atleast one time
\end{quote}

\begin{quote}
\textbf{Jong Hyun (response):} Yes, it's expanded in the introduction.
\end{quote}

\begin{itemize}
\tightlist
\item
  \textbf{Decision:} DFT is already expanded at its very first occurrence in Section 1: ``density functional theory (DFT)\textasciitilde{}\cite{kohn1965dft}''. The convention is to expand at the first paper-wide mention and use the acronym thereafter, which we follow. No paper edit needed.
\item
  \textbf{Daniel's note:} Confirmed on 2026-04-24 --- already expanded before.
\item
  \textbf{Status:} {\color{green!60!black}\faCheckCircle} Already satisfied (no change required).
\end{itemize}

\begin{center}\rule{0.5\linewidth}{0.5pt}\end{center}

\subsection{Section 3 --- Requirements}\label{section-3-requirements}

\subsubsection{3.1 Pranav: optimizer settings and data weights}\label{pranav-optimizer-settings-and-data-weights}

\begin{quote}
\textbf{Pranav:} do we consider optimizer setting here. Setting about data weight
\end{quote}

\begin{itemize}
\tightlist
\item
  \textbf{Decision:} After discussion with Daniel, we extended the ontology with three non-disjoint subclasses of \texttt{Hyperparameter}: \texttt{PhysicalHyperparameter} (direct physical meaning), \texttt{ArchitecturalHyperparameter} (model-architecture choice), and \texttt{TrainingHyperparameter} (training-procedure parameters including optimizer settings and data weights). The non-disjointness is a deliberate modelling choice: many MLIP hyperparameters straddle categories (e.g., \texttt{numRadialBasis} and \texttt{levMax} are both architectural and physical). A strict partition would misrepresent how materials scientists actually think about these parameters.
\item
  \textbf{Daniel's note:} Accepted on 2026-04-24 --- subclasses added with explicit non-disjointness.
\item
  \textbf{Status:} {\color{green!60!black}\faCheckCircle} \textbf{Applied} (commit \texttt{a89570a}).

  \begin{itemize}
  \tightlist
  \item
    XHTML source extended with three Class sections.
  \item
    Running example gained a ``Classifying hyperparameters'' paragraph and example Hyperparameter instances now carry the subclass types.
  \item
    Regenerated \texttt{mlips.ttl} (23 classes) and Appendix B.
  \item
    HermiT consistency verified.
  \end{itemize}
\end{itemize}

\subsubsection{3.2 Xi: ambiguity of ``number of configurations'' (CQ5)}\label{xi-ambiguity-of-number-of-configurations-cq5}

\begin{quote}
\textbf{Xi:} Is this assigned for only one item? e.g., if one says, 100 configurations in the training dataset, it is not clear whether this is along ``chemical'' or ``vibrational'' degree of freedom or coupled by both.
\end{quote}

\begin{itemize}
\tightlist
\item
  \textbf{Decision (postponed):} Moved to end of the list for later discussion. Two options under consideration:

  \begin{itemize}
  \tightlist
  \item
    \textbf{Option A (minimal --- documentation):} Clarify in prose that a ``configuration'' is a snapshot (atomic positions + cell + species) and leave sampling-strategy annotation to future work. Rephrase CQ5 to mention sampling strategy as a question.
  \item
    \textbf{Option B (schema extension):} Add a new property \texttt{samplingStrategy} on \texttt{TrainingDataset} with values like \texttt{Chemical}, \texttt{Vibrational}, \texttt{Mixed}, and optionally split \texttt{numConfigurations} into \texttt{numChemicalConfigurations} and \texttt{numVibrationalConfigurations}.
  \end{itemize}
\item
  \textbf{Daniel's note (2026-04-24):} Not decided yet; revisit at the end.
\item
  \textbf{Resolution (2026-04-25):} Hybrid of Options A and B. The schema is extended now (Option B) with a single optional, multi-valued property \texttt{samplingStrategy:\ TrainingDataset\ $\to$\ SamplingStrategy}; \texttt{SamplingStrategy} is an open extension point (concrete strategies modeled as instances with \texttt{rdfs:label}/\texttt{rdfs:comment}, no committed taxonomy in the ISWC version). CQ5 is rephrased to ask about the sampling strategy (Option A's prose change). Splitting \texttt{numConfigurations} is not done --- the per-strategy counts can come later if needed. The taxonomy of strategies (\texttt{Chemical}, \texttt{Vibrational}, \texttt{Mixed}, \ldots) will be developed in the journal extension.
\item
  \textbf{Status:} {\color{green!60!black}\faCheckCircle} \textbf{Applied} (commit \texttt{8bf888c}).
\end{itemize}

\subsubsection{3.3 Xi: ``material property'' wording in CQ6}\label{xi-material-property-wording-in-cq6}

\begin{quote}
\textbf{Xi:} I feel it is better to say ``on a given material property'' since no studies can cover everything of a given material.
\end{quote}

\begin{itemize}
\tightlist
\item
  \textbf{Decision:} Accept. Benchmarks are always against specific properties (energy, force RMSE, stress, elastic constants, etc.), not ``the material as a whole''.
\item
  \textbf{Daniel's note:} Accepted on 2026-04-24.
\item
  \textbf{Status:} {\color{green!60!black}\faCheckCircle} \textbf{Applied} (commit \texttt{5e25b04}). CQ6 now reads: ``benchmarked a given algorithm on a given material \textbf{property}, and what accuracy metrics were reported?''
\end{itemize}

\subsubsection{3.4 Xi: negative results in CQ7}\label{xi-negative-results-in-cq7}

\begin{quote}
\textbf{Xi:} ``which perform not well'' is also frequently shown in the literature, i.e., convergence test or comparison of different MLIPs. These should be also useful and be answered by the ontology.
\end{quote}

\begin{itemize}
\tightlist
\item
  \textbf{Decision:} Rephrase CQ7 as an explicit \textbf{ranking} query, which cleanly generalizes best, worst, and full-ranking interpretations and removes ambiguity. Best/worst are trivially derivable (reverse the \texttt{ORDER\ BY} direction). Xi's concern is addressed because negative results are now part of the specified capability.
\item
  \textbf{Daniel's note (2026-04-24):} Accepted --- ranking is unambiguous and subsumes all subqueries.
\item
  \textbf{Status:} {\color{green!60!black}\faCheckCircle} \textbf{Applied} (commit \texttt{38eed2f}). Updated in three places:

  \begin{itemize}
  \tightlist
  \item
    Section 3 (CQ7 definition): ``how do different algorithm and hyperparameter combinations rank by accuracy?''
  \item
    Section 8 (Evaluation CQ7 paragraph): heading and question rephrased.
  \item
    Comparison table row: ``Accuracy ranking''.
  \item
    Appendix E (DL queries): subsection heading updated.
  \end{itemize}
\end{itemize}

\subsubsection{3.5 Xi: new CQ on efficiency}\label{xi-new-cq-on-efficiency}

\begin{quote}
\textbf{Xi:} Maybe one more CQ, the efficiency of the algorithms (training cost and running cost). This is also critical for MLIP selection on top of their accuracies.
\end{quote}

\begin{itemize}
\tightlist
\item
  \textbf{Decision:} Accepted with a refinement: Daniel emphasized two distinct dimensions, (a) \textbf{prior knowledge about algorithms} (asymptotic complexity, GPU/parallel support --- properties of the algorithm itself) and (b) \textbf{actual measured costs} (wall-clock time, memory, hardware used --- properties of a specific training run or benchmark result). Both are modelled explicitly.
\item
  \textbf{Daniel's note (2026-04-24):} Accepted --- cover both prior-knowledge and empirical dimensions.
\item
  \textbf{Status:} {\color{green!60!black}\faCheckCircle} \textbf{Applied} (commit \texttt{1eff007}). 10 new datatype properties added:

  \begin{itemize}
  \tightlist
  \item
    On \texttt{MLIPAlgorithm} (prior knowledge): \texttt{trainingComplexity}, \texttt{inferenceComplexity}, \texttt{supportsGPU}, \texttt{supportsParallelization}.
  \item
    On \texttt{TrainingRun} (measured): \texttt{trainingDuration}, \texttt{trainingHardware}, \texttt{peakMemory}, \texttt{gpuHours}.
  \item
    On \texttt{BenchmarkResult} (measured): \texttt{inferenceTimePerAtom}, \texttt{inferenceHardware}.
  \item
    CQ9 added to Section 3 with two-dimensional phrasing; CQ9 SPARQL query added to Appendix C; short explanation added to Section 4 Algorithm Module Core-properties paragraph. HermiT-consistent.
  \end{itemize}
\end{itemize}

\begin{center}\rule{0.5\linewidth}{0.5pt}\end{center}

\subsection{Section 4 --- The MLIPs Ontology}\label{section-4-the-mlips-ontology}

\subsubsection{4.1 Xi: ``MLIP algorithm family'' meaning}\label{xi-mlip-algorithm-family-meaning}

\begin{quote}
\textbf{Xi:} Is here ``MLIP algorithm family'' meant to be each specific ``MLIP'' or a bigger category? I understand ``family'' may also indicate ``common physical background'', i.e., several MLIPs can be grouped into the same physical approximation, e.g., local or non-local, linear or non-linear.
\end{quote}

\begin{itemize}
\tightlist
\item
  \textbf{Decision:} Drop the word ``family''. Daniel also noted that family boundaries are inherently fuzzy --- hybrid algorithms (e.g., mixing moment tensors with message passing) don't fit cleanly into any single family. Trying to encode families rigidly would be a small research project in itself; the paper clarifies granularity in prose instead and defers the classification to future work.
\item
  \textbf{Daniel's note (2026-04-24):} Accepted --- drop ``family''.
\item
  \textbf{Status:} {\color{green!60!black}\faCheckCircle} \textbf{Applied} (commit \texttt{3b91f21}). Changes:

  \begin{itemize}
  \tightlist
  \item
    Section 4 intro: ``models MLIP algorithms'' (not ``algorithm families''); also mentions CQ9.
  \item
    Section 4 Core classes: ``\(\MLIPAlgorithmC\), representing a specific MLIP algorithm (e.g., MACE, ACE, HDNNP)---the architecture itself, not a trained model.''
  \end{itemize}
\end{itemize}

\subsubsection{4.2 Yongliang: fingerprinting the atomic environment}\label{yongliang-fingerprinting-the-atomic-environment}

\begin{quote}
\textbf{Yongliang:} We should mention different MLIPs have different algorithms to fingerprint the atomic environment.
\end{quote}

\begin{itemize}
\tightlist
\item
  \textbf{Decision:} Accepted with a schema extension. Rather than only mentioning descriptors in prose, we introduced a first-class representation: the class \texttt{AtomicEnvironmentDescriptor} and the object property \texttt{hasDescriptor} (domain \texttt{MLIPAlgorithm}, range \texttt{AtomicEnvironmentDescriptor}). Five named individuals populate the initial vocabulary: \texttt{MomentTensorDescriptor}, \texttt{SOAPDescriptor}, \texttt{SymmetryFunctionDescriptor}, \texttt{EquivariantMessagePassingDescriptor}, \texttt{AtomicClusterExpansionDescriptor}. This enables queries like ``which algorithms share the same descriptor?''.
\item
  \textbf{Daniel's note (2026-04-24):} Accepted --- class + property with named individuals for the main descriptor families.
\item
  \textbf{Status:} {\color{green!60!black}\faCheckCircle} \textbf{Applied} (commit \texttt{9ee1ecf}). XHTML source extended; regenerated \texttt{mlips.ttl} (24 classes, 30 object properties); Section 4 Algorithm Module gained a sentence introducing descriptors. HermiT consistency verified.
\end{itemize}

\subsubsection{4.3 Pranav: HDNNP redundancy check}\label{pranav-hdnnp-redundancy-check}

\begin{quote}
\textbf{Pranav:} HDNNP appear here
\end{quote}

\begin{itemize}
\tightlist
\item
  \textbf{Decision:} Pranav flagged that HDNNP appears again in Section 4 after being introduced in Section 1 --- a note-to-self rather than a change request. Standard convention is to introduce an acronym once and use it freely thereafter.
\item
  \textbf{Daniel's note (2026-04-24):} No paper edit needed.
\item
  \textbf{Status:} {\color{green!60!black}\faCheckCircle} No action needed (convention already followed).
\end{itemize}

\subsubsection{4.4 Xi: ``model'' vs ``trained model''}\label{xi-model-vs-trained-model}

\begin{quote}
\textbf{Xi:} maybe ``the trained model''? ``model'' can also be referred to ``MLIP model'' (algorithm).
\end{quote}

\begin{itemize}
\tightlist
\item
  \textbf{Decision:} Accept. The word ``model'' is ambiguous in this domain: the algorithm itself can be called a model (as in ``MTP model''), and the trained artifact is also a model. We consistently use \texttt{TrainedModel} for the artifact; the prose should follow.
\item
  \textbf{Daniel's note (2026-04-24):} Accepted.
\item
  \textbf{Status:} {\color{green!60!black}\faCheckCircle} \textbf{Applied} (commit \texttt{d5147f7}). Sentence now reads ``\ldots the run mediates between algorithm, data, and trained model.''
\end{itemize}

\subsubsection{4.5 Yuji: ``algorithm'' vs ``formalism''}\label{yuji-algorithm-vs-formalism}

\begin{quote}
\textbf{Yuji:} Be aware that the MLIP ``algorithm'' (or rather ``formalism'') does not necessarily fit to the general ML context, because this includes ``how we represent the local atomic environment''.
\end{quote}

\begin{itemize}
\tightlist
\item
  \textbf{Decision:} Very important terminological point. In ML-Schema, \texttt{mls:Algorithm} refers to a training algorithm (e.g., SGD, backpropagation). In the MLIP community, the ``algorithm'' includes the functional form (descriptor + regression scheme), which is more like a ``model formalism'' than a pure training algorithm.
\item
  \textbf{Resolution (2026-04-25):} Adopted a structural fix rather than rewording the alignment paragraph. The class was renamed from \texttt{MLIPAlgorithm} to \texttt{MLIPMethod} and declared \textbf{disjoint} from \texttt{mls:Algorithm} (axiom L11/A11). Three component roles capture what an MLIP method actually is: \texttt{hasFunctionalForm}, \texttt{hasLossFunction}, and \texttt{hasTrainingAlgorithm}, where the last is the \emph{single} alignment point with \texttt{mls:Algorithm}. The same disjointness is why \texttt{MLIPRun} is a subclass of \texttt{prov:Activity} only and not of \texttt{mls:Run}. A new appendix (\texttt{appendix-method-vs-algorithm.tex}) discusses the design.
\item
  \textbf{Status:} {\color{green!60!black}\faCheckCircle} \textbf{Applied} (commit \texttt{5f67fc8}). \texttt{MLIPMethod} disjoint from \texttt{mls:Algorithm}; \texttt{hasTrainingAlgorithm} is the alignment point; \texttt{FunctionalForm} and \texttt{LossFunction} are first-class component classes.
\end{itemize}

\subsubsection{4.6 Pranav: where does PINN fit?}\label{pranav-where-does-pinn-fit}

\begin{quote}
\textbf{Pranav:} Where we will put PINN its input is hyperparameter and output is again hyperparameter for semi-empirical Potential
\end{quote}

\begin{itemize}
\tightlist
\item
  \textbf{Decision:} Physics-informed neural networks for semi-empirical potentials are an interesting edge case where one algorithm produces the hyperparameters of another. The ontology can represent this: a PINN is an \texttt{MLIPAlgorithm}, its output \texttt{TrainedModel} can be used as input (via a new role like \texttt{providesHyperparametersFor}) to another training run. Needs a schema extension.
\item
  \textbf{Resolution (2026-04-25):} Moved to the journal extension (\texttt{2027-swj-mlips-onto}). Detailed placeholder section added at \texttt{sections/11-future-work.tex} §1 (Meta-Learning Patterns), with a draft plan covering three options: (1) provenance-only \texttt{providesHyperparametersFor} property, (2) \texttt{sourcedFrom} on \texttt{HyperparameterSetting} aliased to \texttt{prov:wasDerivedFrom}, or

  \begin{enumerate}
  \def\labelenumi{(\arabic{enumi})}
  \setcounter{enumi}{2}
  \tightlist
  \item
    \texttt{CompositeMethod} subclass with \texttt{bindsHyperparametersOf}. The journal version will pick one mechanism and work out a concrete example with input from Pranav.
  \end{enumerate}
\item
  \textbf{Why not also flag it in the ISWC conclusions:} We considered adding a brief ``future work'' mention of PINN/meta-learning to the ISWC conclusions but decided against it. ISWC Resources Track reviewers are mostly from the semantic-web/ontology community and are unlikely to recognise the PINN-parameterising-a-semi-empirical-potential pattern; flagging the gap there would invite criticism without pre-empting it, whereas the journal version (read by a more materials-aware audience) is the right place to develop the case. Reply to Pranav: comment reviewed, plan recorded, deferred to SWJ extension --- no change to the ISWC paper.
\item
  \textbf{Status:} {\color{green!60!black}\faCheckCircle} \textbf{Deferred to journal extension} --- placeholder and draft plan in place; will be developed as part of the SWJ manuscript.
\end{itemize}

\subsubsection{4.7 Jong Hyun: non-DFT training data}\label{jong-hyun-non-dft-training-data}

\begin{quote}
\textbf{Jong Hyun:} Some MLIPs use data obtained using other than DFT, for example wave function methods such as CCSD(T) as done by Yuji.
\end{quote}

\begin{itemize}
\tightlist
\item
  \textbf{Decision:} Very important. The ontology currently assumes DFT as the ground truth (class \texttt{DFTCalculation}, \texttt{DFTSettings}). We should generalize to \texttt{ReferenceCalculation} with \texttt{DFTCalculation} and \texttt{WaveFunctionCalculation} (CCSD(T), MP2, etc.) as subclasses.
\item
  \textbf{Proposed response:} Agreed. We will generalize \texttt{DFTCalculation} to a supertype \texttt{ReferenceCalculation} and make \texttt{DFTCalculation} a subclass, with siblings for wave-function methods. This is a non-breaking change (SPARQL queries on the supertype still work).
\item
  \textbf{Status:} {\color{orange!80!black}\faPauseCircle} Deferred --- schema change, needs coordination.
\end{itemize}

\subsubsection{4.8 Jong Hyun: XC functional citations}\label{jong-hyun-xc-functional-citations}

\begin{quote}
\textbf{Jong Hyun:} The citations of the xc functionals were added.
\end{quote}

\begin{itemize}
\tightlist
\item
  \textbf{Decision:} Good. PBE (Perdew), LDA (Ceperley), SCAN (Sun) are now cited.
\item
  \textbf{Status:} {\color{green!60!black}\faCheckCircle} Will apply Jong's edit (citations \texttt{perdew1996}, \texttt{ceperley1980}, \texttt{sun2015strongly} added to the running example text).
\end{itemize}

\subsubsection{4.9 Jong Hyun: PAW citation}\label{jong-hyun-paw-citation}

\begin{quote}
\textbf{Jong Hyun:} PAW's citation was added in the caption of Fig. 5.
\end{quote}

\begin{itemize}
\tightlist
\item
  \textbf{Decision:} Good. Blöchl's PAW paper (\texttt{blochl94}) is cited.
\item
  \textbf{Status:} {\color{green!60!black}\faCheckCircle} Will apply Jong's edit (citation added).
\end{itemize}

\subsubsection{4.10 Jong Hyun: Smith 2025 example}\label{jong-hyun-smith-2025-example}

\begin{quote}
\textbf{Jong Hyun:} Is this citation by Smith in the caption of Fig. 7 an example? (This comment command works out of a figure environment.)
\end{quote}

\begin{itemize}
\tightlist
\item
  \textbf{Decision:} Yes, ``Smith et al.~2025'' is a fictional placeholder. The reviewer flagged this: we need a real Ti-Al MLIP benchmark with DOI.
\item
  \textbf{Proposed response:} Thanks for catching this --- it is fictional. We will replace it with a real Ti-Al benchmark reference before submission. Candidates: the Meta-Learn project's own benchmarks, or a published MTP/MACE Ti-Al paper.
\item
  \textbf{Resolution (2026-04-26):} Replaced with Kumar, Körmann, Grabowski, Ikeda. \emph{Machine Learning Potentials for Hydrogen Absorption in TiCr\(_2\) Laves Phases.} Acta Materialia 297, 121319 (2025). doi:10.1016/j.actamat.2025.121319. This paper is co-authored by four authors of the present ontology paper and uses MTP via MLIP-2 with VASP+PAW+PBE --- the same algorithm and DFT toolchain as our running example, so the swap preserves the pedagogical structure. The running example is now TiCr\(_2\)-H (C14, hexagonal Laves phase) instead of Ti-Al; the atom set is \(\{\)Ti, Cr, H\(\}\), energy cutoff 400 eV (matches the paper), 3,766 configurations from the initial active-learning step, RMSE 2.81 meV/atom (the paper's C14 result). Note: no published Ti-Al MTP paper from the IMW Stuttgart group exists; external Ti-Al MTP papers (Qi et al.~2023, Nitol et al.

  \begin{enumerate}
  \def\labelenumi{\arabic{enumi})}
  \setcounter{enumi}{2024}
  \tightlist
  \item
    were rejected in favour of ``their own'' paper.
  \end{enumerate}
\item
  \textbf{Status:} {\color{green!60!black}\faCheckCircle} \textbf{Applied} (commit \texttt{b9a97fc}).
\end{itemize}

\begin{center}\rule{0.5\linewidth}{0.5pt}\end{center}

\subsection{Section 5 --- Knowledge Graph (no longer in the main paper)}\label{section-5-knowledge-graph-no-longer-in-the-main-paper}

The Knowledge Graph section was moved to the appendix because the paper now focuses on the ontology only. The following comments still apply if the KG content is added back later.

\subsubsection{5.1 Pranav: separate DFT codes and workflow codes}\label{pranav-separate-dft-codes-and-workflow-codes}

\begin{quote}
\textbf{Pranav:} It is confusing here, can we separate DFT codes and workflow codes
\end{quote}

\begin{itemize}
\tightlist
\item
  \textbf{Decision:} Good categorization. DFT codes (VASP, GPAW, Quantum ESPRESSO) are first-principles calculators; workflow codes (ASE, LAMMPS) are orchestrators. These should be distinct \texttt{Library} subtypes.
\item
  \textbf{Status:} {\color{orange!80!black}\faPauseCircle} Not applicable in the main paper (KG section is in appendix). Noted for the KG-focused companion paper or future extension.
\end{itemize}

\subsubsection{5.2 Jong Hyun: expand GPAW and other abbreviations?}\label{jong-hyun-expand-gpaw-and-other-abbreviations}

\begin{quote}
\textbf{Jong Hyun:} Do we need to expand the abbreviations here or cite the articles?
\end{quote}

\begin{itemize}
\tightlist
\item
  \textbf{Decision:} Cite them. All the mentioned codes (GPAW, LAMMPS, VASP) have primary papers that should be cited at first use in the KG content.
\item
  \textbf{Status:} {\color{orange!80!black}\faPauseCircle} Not applicable in the main paper. Noted.
\end{itemize}

\subsubsection{5.3 Jong Hyun: public seven work packages?}\label{jong-hyun-public-seven-work-packages}

\begin{quote}
\textbf{Jong Hyun:} Are these seven work packages open to the public?
\end{quote}

\begin{itemize}
\tightlist
\item
  \textbf{Decision:} The ERC project's work package structure is internal; the paper should describe the expert consultation in generic terms (``with domain experts across materials science and ML'') rather than referring to specific internal structure.
\item
  \textbf{Status:} {\color{orange!80!black}\faPauseCircle} Not applicable in the main paper.
\end{itemize}

\subsubsection{5.4 Jong Hyun: mentioning ``the team''}\label{jong-hyun-mentioning-the-team}

\begin{quote}
\textbf{Jong Hyun:} Is it customary in the IT field to mention a specific team in an article? In the materials science field people mention specific topics.
\end{quote}

\begin{itemize}
\tightlist
\item
  \textbf{Decision:} Materials science convention should prevail for this audience. When the KG section returns, we will describe what the team does in terms of topics (MLIP development, materials discovery), not the team itself.
\item
  \textbf{Status:} {\color{orange!80!black}\faPauseCircle} Not applicable in the main paper.
\end{itemize}

\begin{center}\rule{0.5\linewidth}{0.5pt}\end{center}

\subsection{Section 9 --- Sustainability}\label{section-9-sustainability}

\subsubsection{9.1 Yuji: public GitHub from the beginning}\label{yuji-public-github-from-the-beginning}

\begin{quote}
\textbf{Yuji:} Can we from the beginning put the main source code on ``public'' GitHub? The internal Git server is not visible from outside, and I can easily imagine that after a few years the repo will not be maintained.
\end{quote}

\begin{itemize}
\tightlist
\item
  \textbf{Decision:} Strong agreement. Reviewers will also expect a public, accessible repository. The University of Stuttgart's internal Git server is not visible externally, which fails the Resources Track expectation.
\item
  \textbf{Proposed response:} Agreed. The primary source will be on public GitHub (or another public forge like GitLab/Codeberg); the Stuttgart server will be an institutional mirror, not the primary home. We will update Section 9 accordingly.
\item
  \textbf{Status:} {\color{orange!80!black}\faPauseCircle} \textbf{Pending} (decision agreed; the actual prose change in \texttt{sections/09-sustainability.tex} has not yet been made --- current Section 9 makes no mention of GitHub).
\end{itemize}

\begin{center}\rule{0.5\linewidth}{0.5pt}\end{center}

\subsection{Summary}\label{summary}

\begin{longtable}[]{@{}
  >{\raggedright\arraybackslash}p{(\columnwidth - 6\tabcolsep) * \real{0.0714}}
  >{\raggedright\arraybackslash}p{(\columnwidth - 6\tabcolsep) * \real{0.1786}}
  >{\raggedright\arraybackslash}p{(\columnwidth - 6\tabcolsep) * \real{0.1190}}
  >{\raggedright\arraybackslash}p{(\columnwidth - 6\tabcolsep) * \real{0.6310}}@{}}
\toprule\noalign{}
\begin{minipage}[b]{\linewidth}\raggedright
\#
\end{minipage} & \begin{minipage}[b]{\linewidth}\raggedright
Author
\end{minipage} & \begin{minipage}[b]{\linewidth}\raggedright
Section
\end{minipage} & \begin{minipage}[b]{\linewidth}\raggedright
Status
\end{minipage} \\
\midrule\noalign{}
\endhead
\bottomrule\noalign{}
\endlastfoot
1.1 & Pranav + Jong & Intro & {\color{green!60!black}\faCheckCircle} Accept algorithm reordering \\
1.2 & Xi + Jong & Intro & {\color{green!60!black}\faCheckCircle} Expand OWL in prose \\
1.3 & Pranav + Jong & Intro & {\color{green!60!black}\faCheckCircle} Expand RDF in prose \\
2.1 & Yuji + Jong & Related & {\color{green!60!black}\faCheckCircle} Add FAIR citations \\
2.2 & Pranav + Jong & Related & {\color{green!60!black}\faCheckCircle} Already done (DFT expanded in §1) \\
3.1 & Pranav & Reqs & {\color{green!60!black}\faCheckCircle} Hyperparameter subclasses (Phys/Arch/Training) \\
3.2 & Xi & Reqs & {\color{green!60!black}\faCheckCircle} samplingStrategy + SamplingStrategy class added \\
3.3 & Xi & Reqs & {\color{green!60!black}\faCheckCircle} Rephrase CQ6 to ``material property'' \\
3.4 & Xi & Reqs & {\color{green!60!black}\faCheckCircle} Add negative-result clarification to CQ7 \\
3.5 & Xi & Reqs & {\color{green!60!black}\faCheckCircle} Add CQ9 on efficiency \\
4.1 & Xi & Ontology & {\color{green!60!black}\faCheckCircle} Drop ``family''; granularity now in §4 prose \\
4.2 & Yongliang & Ontology & {\color{green!60!black}\faCheckCircle} AtomicEnvironmentDescriptor + hasDescriptor \\
4.3 & Pranav & Ontology & {\color{green!60!black}\faCheckCircle} No action (convention already followed) \\
4.4 & Xi & Ontology & {\color{green!60!black}\faCheckCircle} ``trained model'' disambiguation \\
4.5 & Yuji & Ontology & {\color{green!60!black}\faCheckCircle} MLIPMethod disjoint from mls:Algorithm \\
4.6 & Pranav & Ontology & {\color{orange!80!black}\faPauseCircle} PINN / meta-learning --- future work \\
4.7 & Jong & Ontology & {\color{orange!80!black}\faPauseCircle} Generalize DFTCalculation to ReferenceCalculation \\
4.8 & Jong & Ontology & {\color{green!60!black}\faCheckCircle} XC functional citations \\
4.9 & Jong & Ontology & {\color{green!60!black}\faCheckCircle} PAW citation \\
4.10 & Jong & Ontology & {\color{green!60!black}\faCheckCircle} Replaced Smith 2025 with Kumar et al.~2025 \\
5.1 & Pranav & KG & {\color{red!70!black}\faTimesCircle} Not applicable (KG moved to appendix) \\
5.2 & Jong & KG & {\color{red!70!black}\faTimesCircle} Not applicable (KG moved to appendix) \\
5.3 & Jong & KG & {\color{red!70!black}\faTimesCircle} Not applicable (KG moved to appendix) \\
5.4 & Jong & KG & {\color{red!70!black}\faTimesCircle} Not applicable (KG moved to appendix) \\
9.1 & Yuji & Sustain & {\color{orange!80!black}\faPauseCircle} Public GitHub policy --- agreed, prose not yet done \\
\end{longtable}

\textbf{Totals (25 items):} - {\color{green!60!black}\faCheckCircle} accepted and applied: 19 (including 3.2 schema extension, 4.6 deferral with placeholder, and 4.10 running-example swap) - {\color{orange!80!black}\faPauseCircle} deferred / pending (schema change or pending edit): 2 (4.7, 9.1) - {\color{red!70!black}\faTimesCircle} not applicable: 4 (KG section moved to appendix)

\textbf{Truly pending substantive items (2):}

\begin{longtable}[]{@{}
  >{\raggedright\arraybackslash}p{(\columnwidth - 4\tabcolsep) * \real{0.0732}}
  >{\raggedright\arraybackslash}p{(\columnwidth - 4\tabcolsep) * \real{0.7195}}
  >{\raggedright\arraybackslash}p{(\columnwidth - 4\tabcolsep) * \real{0.2073}}@{}}
\toprule\noalign{}
\begin{minipage}[b]{\linewidth}\raggedright
\#
\end{minipage} & \begin{minipage}[b]{\linewidth}\raggedright
Topic
\end{minipage} & \begin{minipage}[b]{\linewidth}\raggedright
Type
\end{minipage} \\
\midrule\noalign{}
\endhead
\bottomrule\noalign{}
\endlastfoot
4.7 & Generalise \texttt{DFTCalculation} $\to$ \texttt{ReferenceCalculation} & Schema change \\
9.1 & Switch primary source to public GitHub & Prose edit (§9) \\
\end{longtable}


%% ====================================================================
\subsection{Running-Example Phase Choice: C15 (Cubic) over C14 (Hexagonal)}
\label{sec:appendix-internal-c15-vs-c14}

This is an internal-only note recording why the running example uses
the C15 (cubic) phase of TiCr$_2$ rather than the C14 (hexagonal)
phase. The trade-off is small but worth documenting because the
choice cascades through several artefacts (\S\ref{sec:running-example},
Fig.~\ref{fig:simulation-cell}, Fig.~5, Fig.~7, the axiom examples in
Appendix~\ref{sec:axiom-catalog}, and Appendix~\ref{sec:appendix-example-kumar2025}).

\paragraph{Initial choice: C14.}
The Smith~2025-to-Kumar~2025 swap (commit \texttt{b9a97fc}) initially
used the C14 (hexagonal) phase because Kumar et al.'s C14 step-1
training set has 3{,}766 configurations versus C15's 1{,}019, and the
C14-MTP achieves a lower energy RMSE (2.81 meV/atom vs C15-MTP's
3.17 meV/atom). The larger dataset and slightly better accuracy made
C14 look like the more substantial example.

\paragraph{Why we switched to C15.}
Fig.~\ref{fig:simulation-cell-3d} draws a \emph{cubic} simulation cell
(a generic atoms-in-a-box schematic, not a literal Laves cell), but
the prose claimed C14 (hexagonal). For an audience unfamiliar with
the Strukturbericht designations, a cubic-looking figure paired with
the word ``hexagonal'' in the surrounding text is jarring and
distracts from the pedagogical point of the example.

\paragraph{Decision.}
Switch the example to the C15 (cubic) phase (commit \texttt{2352e8a}).
The cubic schematic now matches the prose, and the audience-friendly
opening paragraph calls C15 simply ``the cubic form'' so a reader does
not need to know what ``C15'' or ``Strukturbericht'' means to follow
the example. The cost is the slightly less impressive numbers
(1{,}019 instead of 3{,}766 configs; 3.17 instead of 2.81 meV/atom),
but the goal of \S\ref{sec:running-example} is to make the ontology
encoding clear, not to showcase the most accurate model.

\paragraph{Trade-offs noted, in case we revisit.}
\begin{itemize}
\item \textbf{Pedagogical clarity (chosen).} C15 cubic matches the
  cubic schematic in Fig.~\ref{fig:simulation-cell-3d}; an
  ISWC/SemWeb reader does not need crystallography to follow the
  example.
\item \textbf{Dataset substantiveness (rejected).} C14 is the larger
  training set, and the C14-MTP is the more accurate model. If we
  ever redraw Fig.~\ref{fig:simulation-cell-3d} as a faithful C14
  hexagonal cell, switching back to C14 would reclaim that.
\item \textbf{Phase as a first-class concept.} Either way, the C15
  vs.\ C14 distinction lives in a \texttt{materialClass} string,
  not in a structured \dlConcept{Phase} concept. Adding a
  CMSO/EMMO-aligned crystal-structure slot is on the journal-extension
  list (\S\ref{sec:appendix-example-kumar2025-missing}).
\end{itemize}

%% ====================================================================
